\vspace{-2mm}

\begin{table}[t]
\centering
\resizebox{\ifdim\width>\textwidth\textwidth\else\width\fi}{!}{

\begin{tabular}{cccclcccccccccc}
\toprule[1.2pt]
& \multicolumn{3}{c}{Methods} & & \\
\cmidrule[0.6pt]{2-4}
& \multirow{2}{*}{\begin{tabular}[c]{@{}c@{}}Non-linear \\ message passing\end{tabular}} & \multirow{2}{*}{\begin{tabular}[c]{@{}c@{}}MLP \\ attention\end{tabular}} & \multirow{2}{*}{\begin{tabular}[c]{@{}c@{}}Dot product \\ attention\end{tabular}} &  Task & $\alpha$ & $\Delta \varepsilon$ & $\varepsilon_{\text{HOMO}}$ & $\varepsilon_{\text{LUMO}}$ & $\mu$ & $C_{\nu}$ & & Training time & Number of \\
%Index & & & & Unit & bohr$^3$ & meV & meV & meV & D & cal/mol K & & (minutes/epoch) & parameters \\
Index & & & & Unit & $a_0^3$ & meV & meV & meV & D & cal/mol K & & (minutes/epoch) & parameters \\
\midrule[1.2pt]

1 & \ding{51} & \ding{51} &  & & .046 & 30 & 15 & 14 & .011 & .023 & & 12.1 & 3.53M \\
2 & & \ding{51} &   & & .051 & 32 & 16 & 16 & .013 & .025 & & 7.2 & 3.01M \\
3 & & & \ding{51} & & .053 & 32 & 17 & 16 & .013 & .025 & & 7.8 & 3.35M \\
%4 & & \ding{51} &  &  & & .060 &  & 18 & 17 & .016 &  \\

\bottomrule[1.2pt]

\end{tabular}
}
\vspace{-3mm}
\caption{\textbf{Ablation study results on QM9.}}
\vspace{-3.25mm}
\label{tab:qm9_ablation}
\end{table}


\begin{table}[t]
\begin{adjustwidth}{0cm}{0cm}
\centering
\resizebox{\ifdim\width>\textwidth\textwidth\else\width\fi}{!}{

\begin{tabular}{cccccccccccccccccc}
\toprule[1.2pt]
& \multicolumn{3}{c}{Methods} & & \multicolumn{5}{c}{Energy MAE (eV) $\downarrow$} & & \multicolumn{5}{c}{EwT (\%) $\uparrow$} \\
\cmidrule[0.6pt]{2-4}
\cmidrule[0.6pt]{6-10}
\cmidrule[0.6pt]{12-16}
& \multirow{2}{*}{\begin{tabular}[c]{@{}c@{}}Non-linear \\ message passing\end{tabular}} & \multirow{2}{*}{\begin{tabular}[c]{@{}c@{}}MLP \\ attention\end{tabular}} & \multirow{2}{*}{\begin{tabular}[c]{@{}c@{}}Dot product \\ attention\end{tabular}} & & & & & & & & & & & & & Training time & Number of \\

Index & & & &  & ID & OOD Ads & OOD Cat & OOD Both & Average & & ID & OOD Ads & OOD Cat & OOD Both & Average & (minutes/epoch) & parameters \\
\midrule[1.2pt]

1 & \ding{51} & \ding{51} &  & & 0.5088 & 0.6271 & 0.5051 & 0.5545 &  0.5489  & & 4.88 & 2.93 & 4.92 & 2.98 & 3.93 & 130.8 & 9.12M \\
2 & & \ding{51} &   & & 0.5168 & 0.6308 & 0.5088 & 0.5657 & 0.5555 &  & 4.59 & 2.82 & 4.79 & 3.02 & 3.81 & 91.2 & 7.84M \\
3 & & & \ding{51}  & & 0.5386 & 0.6382 & 0.5297 & 0.5692 & 0.5689 & & 4.37 & 2.60 & 4.36 & 2.86 & 3.55 & 99.3 & 8.72M \\
%4 & & \ding{51} &  &  & 0.5213 & 0.6268 & 0.5178 & 0.5621 & 0.5579\\

%\bottomrule[1.2pt]
%\bottomrule[1.2pt]
%\vspace{-1mm}
%\\
%\toprule[1.2pt]

%%%\midrule[1.2pt]
%%%& \multicolumn{3}{c}{Methods} & & \multicolumn{5}{c}{EwT (\%) $\uparrow$} \\
%%%\cmidrule[0.6pt]{2-4}
%%%\cmidrule[0.6pt]{6-10}
%%%& \multirow{2}{*}{\begin{tabular}[c]{@{}c@{}}Non-linear \\ message passing\end{tabular}} & \multirow{2}{*}{\begin{tabular}[c]{@{}c@{}}MLP \\ attention\end{tabular}} & \multirow{2}{*}{\begin{tabular}[c]{@{}c@{}}Dot product \\ attention\end{tabular}} &  \\

%%%Index & & & & & ID & OOD Ads & OOD Cat & OOD Both & Average \\
%%%\midrule[1.2pt]

%%%1 & \ding{51} & \ding{51} & & & 4.88 & 2.93 & 4.92 & 2.98 & 3.93 \\
%%%2 & & \ding{51} &   &  & 4.59 & 2.82 & 4.79 & 3.02 & 3.81  \\
%%%3 & & & \ding{51} & & 4.37 & 2.60 & 4.36 & 2.86 & 3.55 \\
%4 & & \ding{51} &  &  & 4.37 & 2.80 & 4.40 & 2.98 & 3.64 \\

\bottomrule[1.2pt]

\end{tabular}
}

\vspace{-3mm}
\end{adjustwidth}
\caption{\textbf{Ablation study results on OC20 IS2RE validation set.}
%\note{Dot product attention uses smaller learning rate $1.5 \times 10^{-3}$.}
}
\vspace{-6.75mm}
\label{tab:oc20_ablation}
\end{table}



\vspace{-2.25mm}
\subsection{Ablation Study}
\vspace{-2.5mm} 
\label{subsec:ablation_study}

We conduct ablation studies \revision{to show that Equiformer with dot product attention and linear message passing has already achieved strong empirical results and demonstrate} the improvement brought by MLP attention and non-linear messages in the proposed equivariant graph attention.
%We modify dot product (DP) attention~\citep{transformer, se3_transformer} so that it only differs from MLP attention in how attention weights $a_{ij}$ are generated from $f_{ij}$.
\revision{Dot product (DP) attention only differs from MLP attention in how attention weights $a_{ij}$ are generated from $f_{ij}$.}
%%%Please refer to Sec.~\ref{appendix:subsec:dot_product_attention} in appendix for details on DP attention.
Please refer to Sec.~\ref{appendix:subsec:dot_product_attention} in appendix for further details.
For experiments on QM9 and OC20, unless otherwise stated, we follow the hyper-parameters used in previous experiments. 
%\revision{Please refer to Sec.~\ref{appendix:subsec:qm9_training_details} and Sec.~\ref{appendix:subsec:oc20_training_details} for training time of different models.}

\vspace{-3.25mm}
\paragraph{Result on QM9.}
The comparison is summarized in Table~\ref{tab:qm9_ablation}.
\revision{Compared with models in Table~\ref{tab:qm9_result}, Equiformer with dot product attention and linear message passing (Index 3) achieves competitve results.}
%Non-linear messages improve upon linear messages when MLP attention is used.
Non-linear messages improve upon linear messages when MLP attention is used while non-linear messages increase the number of tensor product operations in each block from $1$ to $2$ and thus inevitably increase training time.
%%%Similar to what is reported by GATv2~\cite{gatv2}, the improvement of replacing DP attention with MLP attention is not very significant.
%%%\todo{On the other hand, the improvement of replacing DP attention with MLP attention is not very significant.}
On the other hand, MLP attention achieves similar results to DP  attention.
We conjecture that DP attention with linear operations is expressive enough to capture common attention patterns as the numbers of nighboring nodes and atom species are much smaller than those in OC20.
However, MLP attention is roughly $8\%$ faster as it directly generates scalar features and attention weights from $f_{ij}$ instead of producing additional key and query irreps features for attention weights.

\vspace{-3.25mm}
\paragraph{Result on OC20.}
%%%We consider the setting of training without IS2RS auxiliary task and use a smaller learning rate $1.5 \times 10 ^{-4}$ for DP attention as this improves the performance.
We consider the setting of training without auxiliary task and summarize the comparison in Table~\ref{tab:oc20_ablation}.
%%%We summarize the comparison in Table~\ref{tab:oc20_ablation}.
\revision{Compared with models in Table~\ref{tab:is2re_validation}, Equiformer with dot product attention and linear message passing (Index 3) has already outperformed all previous models.}
Non-linear messages consistently improve upon linear messages.
In contrast to the results on QM9, MLP attention achieves better performance than DP attention and is $8\%$ faster.
We surmise this is because OC20 contains larger atomistic graphs with more diverse atom species and therefore requires more expressive attention mechanisms.
Note that Equiformer can potentially improve upon previous equivariant Transformers~\citep{se3_transformer, torchmd_net, eqgat} since they use less expressive attention mechanisms similar to Index 3 in Table~\ref{tab:oc20_ablation}.